\chapter{Administrer le réseau de recharge}

\section{Inventaire des stations}

La page \emph{Stations} propose des vues en cartes, liste, tableau et carte
géographique. Elle permet de filtrer les actifs, d'ouvrir leur détail et de
lancer le commissionnement d'une station.

\manualscreenshotcompact{admin-stations.png}
  {Inventaire des stations de l'organisation}
  {fig:admin-stations}

Une station possède une référence métier, une position, des caractéristiques
matérielles, une identité OCPP et un ou plusieurs connecteurs physiques.

\section{Commissionner une station simulée}

Le parcours de commissionnement crée de façon atomique la station, ses
connecteurs et les informations nécessaires à sa connexion OCPP. Dans le cadre
du MVP, le parcours privilégié utilise le simulateur SAP local.

\begin{procedurebox}{Créer une station reliée au simulateur local}
  \begin{enumerate}[leftmargin=*]
    \item Ouvrir \emph{Stations} puis choisir \emph{Add station}.
    \item Renseigner le nom, une référence unique, l'adresse et la position sur
    la carte.
    \item Choisir un profil matériel simulé. Le modèle, la puissance et les
    connecteurs sont définis par ce profil vérifié.
    \item Vérifier que l'identité OCPP est unique. Elle correspond exactement à
    celle utilisée par le simulateur dans l'URL WebSocket.
    \item Vérifier le récapitulatif puis créer la station.
    \item Suivre l'état de provisionnement automatique affiché par l'interface.
    \item Ouvrir le détail lorsque la borne est connectée.
  \end{enumerate}
\end{procedurebox}

\expected{Après le \emph{BootNotification} et les premiers battements de cœur,
la connexion devient active et l'état des connecteurs est projeté à partir des
événements OCPP.}

\attention{La référence \texttt{CT-TUN-001} existe déjà dans le jeu de données
de démonstration. Choisir une nouvelle référence, par exemple
\texttt{CT-TUN-010}. La référence et l'identité OCPP doivent rester uniques.}

\section{Utiliser le laboratoire de simulation OCPP}

Le \emph{Simulation Lab} est un espace de test indépendant des tableaux de
bord métier. Il permet de piloter les stations simulées sans ouvrir un terminal
et de visualiser les signaux qui traversent la passerelle OCPP. La carte des
signaux représente des événements assainis : elle n'expose ni secret de borne,
ni jeton de contrôle, ni charge utile technique confidentielle.

\manualscreenshot{admin-simulation-lab.png}
  {Laboratoire de simulation OCPP et signaux d'une station}
  {fig:admin-simulation-lab}

\begin{table}[H]
  \centering
  \caption{Familles d'actions disponibles dans le laboratoire}
  \label{tab:simulation-lab-actions}
  \begin{tabularx}{\textwidth}{|L{4cm}|Y|Y|}
    \hline
    \rowcolor{ctmint}
    \textbf{Famille} & \textbf{Exemples} & \textbf{Résultat à observer} \\ \hline
    Connexion de la station & Connexion, déconnexion et battement de cœur &
    État OCPP, dernier signal et règle de disponibilité. \\ \hline
    État physique du connecteur & Branchement, débranchement, défaut et
    récupération & Projection du connecteur et disponibilité de la station.
    \\ \hline
    Scénarios reproductibles & Cycle de câble et récupération après défaut &
    Suite ordonnée d'événements dans le flux de protocole. \\ \hline
    Commandes centrales & Redémarrage logiciel, déverrouillage et maintenance &
    Accusé de réception, effet simulé et trace de commande. \\ \hline
  \end{tabularx}
\end{table}

\begin{procedurebox}{Tester un changement d'état simulé}
  \begin{enumerate}[leftmargin=*]
    \item Ouvrir \emph{Simulation Lab} et sélectionner une station de
    l'organisation.
    \item Vérifier que l'adaptateur et le lien OCPP sont opérationnels.
    \item Sélectionner le connecteur physique à tester.
    \item Lancer une action unique ou un scénario reproductible.
    \item Suivre les impulsions, le flux du protocole et l'état projeté.
    \item Revenir au détail de la station pour contrôler la disponibilité, les
    alertes et l'historique de commande.
  \end{enumerate}
\end{procedurebox}

L'administrateur dispose du contrôle des stations simulées de son
organisation. L'opérateur utilise les actions d'exploitation qui lui sont
autorisées ; le technicien limite ses essais aux diagnostics liés à son
intervention. Le client n'accède pas à cette console administrative : lors du
parcours de recharge, l'étape de branchement observe automatiquement l'état
OCPP du connecteur sélectionné.

\attention{Le laboratoire reproduit le comportement attendu d'une borne, mais
ne constitue pas une preuve de compatibilité avec un équipement physique. Un
essai doit porter sur une station simulée identifiée et ne doit pas perturber
une session en cours.}

\section{Connecteurs}

Le connecteur présenté dans l'interface correspond à un connecteur physique de
la station. Son identifiant OCPP doit correspondre au \texttt{connectorId}
envoyé par la borne ou le simulateur.

\begin{table}[H]
  \centering
  \caption{Informations principales d'un connecteur}
  \label{tab:connector-fields}
  \begin{tabularx}{\textwidth}{|L{3.5cm}|Y|Y|}
    \hline
    \rowcolor{ctmint}
    \textbf{Champ} & \textbf{Exemple} & \textbf{Règle} \\ \hline
    Libellé & A1 & Nom lisible par l'utilisateur. \\ \hline
    Identifiant OCPP & 1 & Valeur numérique envoyée par la station. \\ \hline
    Type de prise & CCS2 & Compatibilité physique du véhicule. \\ \hline
    Type de courant & DC & AC ou DC selon l'équipement. \\ \hline
    Puissance maximale & 120 kW & Capacité nominale, non garantie pendant une
    session. \\ \hline
    État & Available & Valeur projetée depuis les événements OCPP. \\ \hline
  \end{tabularx}
\end{table}

Le QR code d'un connecteur représente un lien de démarrage vers la station et
le connecteur concernés. Dans un déploiement physique, l'étiquette est imprimée
et apposée sur l'équipement ; le client la scanne depuis son téléphone.

\section{Disponibilité calculée}

L'état visible n'est pas un simple champ modifié manuellement. Il résulte des
événements reçus et de règles métier, notamment :

\begin{itemize}[leftmargin=*]
  \item \texttt{StatusNotification} détermine l'état d'un connecteur ;
  \item une transaction active rend le connecteur occupé ;
  \item un défaut déclaré peut rendre le connecteur ou la station indisponible ;
  \item l'absence prolongée de battement de cœur rend la station hors ligne ;
  \item un mode maintenance exclut volontairement l'équipement du service.
\end{itemize}

\section{Documents et suppression}

L'onglet \emph{Documents} contient uniquement les fichiers réellement importés
pour la station : fiche technique, certificat, procédure ou preuve
d'intervention. Une liste vide signifie qu'aucun document n'a encore été
joint.

\attention{Avant de supprimer une station ou un connecteur, vérifier les
sessions, alertes, interventions, documents et traces OCPP associés. Préférer
le mode maintenance ou l'archivage lorsqu'un historique doit être conservé.}

\section{Commandes distantes}

Chaque commande possède un historique distinct de son accusé de réception. Une
commande \emph{Accepted} signifie que la borne a accepté la demande, pas
nécessairement que l'effet physique final est déjà terminé.

\begingroup
\small
\renewcommand{\arraystretch}{1.05}
\begin{longtable}{|L{3.7cm}|L{5.2cm}|L{6.2cm}|}
  \hline
  \rowcolor{ctmint}
  \textbf{Commande} & \textbf{Effet recherché} & \textbf{Précaution} \\ \hline
  \endfirsthead
  \hline
  \rowcolor{ctmint}
  \textbf{Commande} & \textbf{Effet recherché} & \textbf{Précaution} \\ \hline
  \endhead
  Restart station &
  Demander un redémarrage logiciel ou matériel selon la capacité de la borne. &
  Vérifier l'absence de session critique et suivre le retour OCPP. \\ \hline
  Unlock connector &
  Demander le déverrouillage du câble d'un connecteur. &
  Utiliser uniquement si le véhicule et la zone sont sécurisés. \\ \hline
  Set maintenance mode &
  Retirer temporairement la station ou le connecteur du service. &
  Documenter le motif et prévoir la remise en service. \\ \hline
  Remote start &
  Demander le démarrage d'une transaction autorisée. &
  Vérifier le client, le connecteur, le tarif et l'autorisation. \\ \hline
  Remote stop &
  Demander l'arrêt d'une transaction active. &
  Confirmer la session ciblée et le motif de l'arrêt. \\ \hline
  Change availability &
  Demander une disponibilité opérationnelle ou une indisponibilité. &
  Le résultat final dépend de l'acceptation par la borne. \\ \hline
\end{longtable}
\endgroup
